\subsection{Installment 11: “Folkebladet” May 22, 1918}

I do not find that the editorial office fully adheres to the view of the Leading Principles held by their authors. The editor’s earlier reply to Birkeland and his two most recent editorial articles, “The Congregation,” show no change in his understanding of the principles—nor any understanding of what I have been striving to show. To be sure, he states that we are agreed on three or four points. But this agreement is, at the very most, an agreement in a formal, not in a real, sense. The editor’s concession that the principles are “essentially constitution” means, for the conclusions he draws, no more than if he had written: the principles are “essentially water.” The same I would say of the points 1 and 2 which he sets forth: they are immaterial to the consequences the editor deduces. The fourth point—“The church’s principal task is nevertheless not to labor for constitution, but for the salvation of souls, for spirit and life”—with this, however, I am heartily in agreement. For that is, or ought to be, the self-evident principal task of all congregations, churches, and church bodies, indeed of every single man and woman who is called by Christ’s name.

Otherwise I see no principial agreement in our understanding of the principles. The editorial office’s understanding of principles no. 1 and 2 is not in harmony with what I have cited from S. and D. on pages 251 and 301 of Folkebladet (second column) and in many other places. Sverdrup never taught that where there is one believing man, there the congregation is present “according to its essence.” Two Christians in one place are—indisputably—not a congregation, either according to essence or according to form. But if they come together around the Word and the Sacraments in order to help one another, then they are a congregation according to the essence; and if they organize the outward congregational organization, then they are a congregation also according to the form. That was Sverdrup’s doctrine and that is the doctrine of the principles.

The editor confuses the inward essence of the kingdom (righteousness, peace, and joy in the Holy Spirit) with the essence of the congregation (the body of Jesus manifested in the world). And he confuses the outward form of the kingdom (the congregation) with the form of the congregation (the outward organization).

In 1875 the faculty contended (to use its own words) “against pietism and false spirituality, which dissolves Christian society into isolated individuals.” It contended for what it regarded as truth, that “it belongs essentially to the essence of the Christian church to exist in the form of congregations, not as individual Christians standing alone.” There is nothing to indicate that the editor has ever understood this language, the main principle in the Augustana understanding of the congregation. Had he done so, then he would have understood both the principles and their authors, both Birkeland and the spirit at the Fargo meeting, better than he now does; nor would he have voted with Dr. Jacobs as he has done.

Now Prof. Nydahl has the right to hold his opinion. No one can challenge him for that. And he has the right to criticize the principles—under his own name. But the matter becomes somewhat different when he gives editorial interpretations of the principles that are interpretations away from them, and when he puts the brakes on Birkeland’s writings and labels them under dictatorship, absolutism, and so forth.

For nine years now I have been a theological teacher at Augsburg, where the theological faculty has charged me with lecturing on various theological subjects—church polity as well. I have felt it to be a duty to teach in accordance with the Leading Principles. It is possible that I do not possess the qualifications to teach this subject, and that the editorial office has been aware of that through all these years. There can now no longer be any doubt that the editorial office thinks it could be done better, and that there are those at Augsburg who could do it better. I hope that the proper persons concerned will go to the proper authority and bring about a better arrangement. What I wish to make clear is that I share Birkeland’s interpretation of the principles, but not that of the editorial office.

In maintaining my position in my teaching I have gone both to Scripture and to the literature on polity and to much else besides, in order to illuminate the individual principles. I know their strength and their weakness. About their weakness I have kept silent. But I have always interpreted them as I believe their authors intended them. Of the complete conformity of the principles to Scripture and of their absolute validity I have never made any assertion.

In my contributions to the paper I have both in season and out of season cited S. and D.—not in order to prove the conformity of the principles to Scripture (for that I do not need the help of others), but to show that in my understanding and interpretation of the principles I have been faithful to the authors. I have not misrepresented them. “Scripture interprets Scripture; author interprets author;” therefore I also published the little collection called “Guidance,” which has been a guide both to me and to many students, and which could have been the same to the editorial office.

If in these years I have not been able to interpret the principles rightly, then I have never deserved to lecture on theology or to teach anything whatever. If I am so spiritually blind, so scientifically unfit, indeed so intellectually incapable, that through all this time I have been unable to provide a correct diagnosis of the principles, then that is a sorrowful commentary on the annual meetings’ filling of professorships, and an even more sorrowful commentary on our whole higher educational system, both domestic and foreign; for I have, after all, tested a little of both.

I mentioned something about the weakness of the principles. It can be risky to suggest anything of that sort; it can be taken as though the disciple would be above the master. But one must reckon with imperfections in masters also. I shall illustrate this. In Sweden there was carried on for more than a generation a sharp conflict between the theological faculty in Lund and the theological faculty in Uppsala. It concerned the concept of the church. We have mentioned a little about it before.

The church-concept of the Lund faculty was high-church; the means of grace constituted the church, the office was given in and with the means of grace, which were thus regarded as instituted by God; the congregation, the believers, were a product of the means of grace and the office.

The church-concept of the Uppsala faculty was low-church: the means of grace were a product of the believing congregation, the office likewise; the congregation chose a man for the office, but the holder of the office was to be a man with a gift of grace.

A massive literature was the result of the collisions between these two concepts of the church. Each of the two theological camps developed and fortified its own concept of the church and pointed out the weaknesses in the concept embraced by the opponent. Able men, far more isolated on the academic field than the students and professors at the Scandinavian-American schools, shook themselves out again and again for attack and defense. And now? Now the dust has settled. It has been shown that neither of these concepts of the church could stand. In each of them there was a hidden dualism.

Is there any dualism in the Leading Principles? If we are like the sort of Americans of whom it is said, The American cares only for the acts of the apostles, and the acts of his grandfather,—then we shall not want to know of any such question. If we are Augustanas in such a way that we care only for the acts of the apostles and the acts of Sverdrup and Oftedal—that is, care little about history and the continuity of history—then it is daring to raise such a question in our midst.

But let anyone take up Luther, Sohm, W. Walther, and he cannot avoid this question. And if he finds the answer—may he be allowed to state it? When it is so difficult to be permitted to express one’s agreement with the principles, as Birkeland does, how much more difficult will it not be to document a possible disagreement? The times to come must answer that.

The principles have their significance. But they are not what is central in Christianity. And neither B. nor I have made them that. To be sure, it may look as though they were the central thing, if one is to judge from all the column-space they have received from B. and me, especially from me. Yet they have always occupied much space in Folkebladet, even before I began to write. Still, this I will emphasize: the Free Church’s Leading Principles “lay the chief weight upon the congregation and its spirit, life, and freedom.” For those who think that constitutional principles concern only outward arrangements, this sentence, taken from Guidance 118, ought to be shown a measure of serious reflection. Individual ethics aims at the individual’s spirit, life, and freedom. Polity aims at the congregation’s spirit, life, and freedom, although—what of course no one will deny—the former is a condition for the latter. But the principles intend more than this condition: they intend a congregational life brought forth by the Spirit and lived in freedom.

One may be saved without knowing the principles. But, as said, they have their significance, but only as a means to an end, which also holds true of the organized church, great or small, of the office, of the public service of worship including preaching and hymn-singing, of the confessional writings, yes, also of Augsburg Seminary, which one has not—rhetorically at least—placed in the periphery, but has outright made into “the Free Church’s heart.”

Compared with Christ, who is what is central in Christianity, everything else is peripheral. Constitution is no exception, and therefore neither are the Leading Principles, nor Augsburg Seminary either—which, moreover, ought to be superfluous to say.

The great question does not become: Does one stand in harmony with this or that group of principles, this or that wing of an aristocracy of office, this or that institution? But the question—the real one—becomes: Does one stand in accord with the future of God’s kingdom, and is one willing to let oneself be used for the kingdom’s advance, to labor for the fulfillment of the prayer: hallowed be your name; your kingdom come; as in heaven so also on earth—your will be done!

(Conclusion.)
